\documentclass{ctexbook}
\usepackage{amsmath,amssymb,amsthm}
\usepackage[hidelinks]{hyperref}
\newcommand{\Spec}{\operatorname{Spec}}
\newcommand{\codim}{\operatorname{codim}}
\newcommand{\prof}{\operatorname{prof}}
\newcommand{\coprof}{\operatorname{coprof}}
\newcommand{\supp}{\operatorname{supp}}
\newcommand{\ednote}[1]{\footnote{英译编者注：#1}}
\newcommand{\tnote}[1]{\par\smallskip\noindent\textbf{本批译校注：}#1\par\smallskip}
\newcommand{\entry}[2]{\par\bigskip\noindent\textbf{#1\quad #2}\quad}
\title{代数几何基础第五卷预备笔记\\第8节（8.1至8.6）独立续译与译校说明}
\author{据用户上传英译编辑稿译出}
\date{2026年11月14日}
\begin{document}
\maketitle
\frontmatter
\chapter*{范围与底本}
本批连续翻译用户上传的 \texttt{EGAVaiello.pdf} 第一篇第8节的（8.1）至（8.6），包含（8.4.1）、证明及英译编者脚注29--34，位置为底本第21--25页。止于（8.6）之后的句子“前述一般判据在（8.6）的情形应当是显然的”；不包括随后以“We now study the points”开头的未编号段落，亦未译至（8.7）。它不是第8节全文，更不是整部EGA的完成本。

底本为 Piotr Blass 与 Joseph Blass 翻译编辑的 Grothendieck 第五卷预备笔记。本批未声称已与法文手稿校勘。正文保留难辨字、残句及未经证实的数学写法；英译编者说明与本批译校注分开。$P^\vee$ 表示超平面参数空间；承前文的 $F,G_\xi,\phi_\xi$ 不在本批重新定义。

本批接在已经独立交付的第7节之后，与最新第五版合订本分支分离。核查该版实际进度文件可知，其第五卷仍为第1--6节，第0章至（18.5.4）。本批不覆盖、改写或自动合入原工程。因底本存在实质缺文和矛盾，本件是忠实译文附审校记录，不是对全部命题正确性的认证。
\mainmatter
\chapter*{§8\quad 例外超平面集合的维数}
\entry{8.1}{}
在前面各节，尤其是第2节和第3节，我们已给出一些断言：使得 $Y$ 具有某个性质 $P$ 的 $\xi\in P^\vee$ 所成的集合是可构造的，且包含泛点 $\eta$；或者说，对性质 $P$ 而言“例外”的 $\xi\in P^\vee$ 所成的集合 $Z_P$ 是可构造且稀疏的，即其闭包的余维至少为一。〔注意：这里假设 $S=\Spec(k)$。〕
\tnote{底本首句重复了“使得……的 $\xi\in P^\vee$ 所成的集合”这段引导语。本译文只排一次，在此记录重复。}
在某些情形，可以给出这个余维的一个更好的上界，从而使该断言更为精确；这对于某些问题很重要。例如，如果看出这个余维大于或等于二，那么 $P^\vee$ 中一条“充分一般”的直线 $D$ 就不与 $Z_P$ 相交，因而存在（若 $k$ 无限）超平面截面 $Y_\xi$ 的“线性束”（$\xi$ 为 $D$ 的几何点），其中所有截面都具有性质 $P$；例子见第〔节号缺失〕节。\ednote{节号遗漏，请询问 A.G.。〔底本脚注29。〕}
\tnote{“上界”按底本 upper bound 译出；下文却强调余维至少二，方向并不协调。未静默改为“下界”。}
从编写的角度看，本节的结果使前几节的结果更精确，因而产生一个问题：是否应当把这些补充另列一节（或一编号），还是在前文逐步推进时就给出更精确的版本？\textit{Redactor decidetur}〔拉丁文〕。\ednote{由编辑决定。〔底本脚注30。〕}

\entry{8.2}{}
令 $Z$ 为满足 $\dim Y_\xi>\dim X-1$ 的 $\xi\in P^\vee$ 所成的集合。假设对 $X$ 的每个不可约分支 $X_i$，都有
\[
\dim f(X_i)>\text{〔字迹不清：是二吗？请询问 A.G.〕}\text{。}
\]
那么 $Z$ 在 $P^\vee$ 中的余维为二。这由（2.1）、（2.2）（后者蕴含 $Y$ 的每个不可约分支都支配 $P^\vee$）以及态射 $Y\to P^\vee$ 的维数理论得出。从这个结果出发，可以把下述情形作为推论：从 $X$ 的一个闭子集 $Z$ 出发，考察其在 $Y_\xi$ 中的逆像 $Z_\xi$ 的维数（$\xi\in P^\vee$）；甚至可以取 $Z$ 为如下 $\xi\in P^\vee$ 的集合：$T_{k(\xi)}$ 存在一个不可约分支，其在 $Y_\xi$ 上的迹具有最大的维数。〔注意：始终假设对 $Z$ 的每个不可约分支 $Z_i$ 有 $\dim f(Z_i)>0$。〕

最后，这个方向上最精确的断言，是如下修改后的断言；它容易由第一个断言（$X$ 不可约的情形）和（2.7）推出：设 $F$ 是 $X$ 上的凝聚模，假设对 $F$ 的每个相伴素循环 $T$ 都有 $\dim f(T)>0$，则使 $\phi_\xi$ 不是 $F_{k(\xi)}$-正则的 $\xi\in P^\vee$ 所成的集合是（可构造的，并且）余维至少为二。（$\phi_\xi$ 的记号来自第5节。）可以将此作为主要断言，把前面的断言列作推论，而证明则由其中一个推论给出或经由它进行。

请注意，沿用上述记号，若 $\xi\in P^\vee-Z$，则对每个 $y\in Y_\xi$ 都有
\[
\coprof_y G_{k(\xi)}=\coprof_y G_\xi\text{。}
\]
因此，若 $\coprof F\leq n$，则对 $\xi\in P^\vee-Z$，有 $T_\xi$ 的余深度 $\leq n$；特别地，若 $F$ 是 Cohen--Macaulay 的，则对 $\xi\in P^\vee-Z$，$G_\xi$ 也是 Cohen--Macaulay 的。最后，若 $F$ 满足 $(S_k)$，则对 $\xi\in P^\vee-Z$，$G_\xi$ 满足 $(S_{k-1})$（参见 $0_{\mathrm{IV}}$）。
\tnote{本段保留 $T_{k(\xi)}$、$G_{k(\xi)}$ 和 $T_\xi$ 等不协调的符号，未把它们猜改为 $F$ 或 $G$。“余维为二”亦是底本原说法，不视为已证明的精确等式；例如恒等嵌入 $X=P=\mathbf P^3$ 的异常集为空。}

\entry{8.3}{}
我们注意到：如果 $F$ 满足 $(S_k)$，对于某个 $\xi\in P^\vee$，使得 $\phi_\xi$ 是 $f_{k(\xi)}$-正则的，并且 $G_\xi$ 有一个余维至少为二的分支\ednote{字迹不清，请询问 A.G.。〔底本脚注31。〕}，即使 $F=\mathcal O_X$、$k=1$，且 $X$ 几何整、维数为二；或者（$k=2$，$X$ 几何整且几何正规、维数为三）。〔此句按底本保留残缺结构。〕
只须从代数闭域 $k$ 上的一张射影整曲面
\[
X\subset P^r
\]
出发，它有一点 $x$，使 $X$ 在该点不是 Cohen--Macaulay 的；那么，对每个经过 $x$ 的超平面，相应的超平面截面 $Y_\xi$ 都以 $x$ 为一个嵌入相伴循环。相应地，从三维正规（因此满足 $(S_2)$）整簇 $X\subset P^\vee$ 出发，它有一点 $X\in X$，使 $X$ 在此处不是 Cohen--Macaulay 的；则经过 $x$ 的那些 $Y$ 不是 CM 的，也就是说，它们在 $x$ 处“met $(S_2)$”〔底本作 they are met $(S_2)$，英文用语不明〕。

在这些例子中，对性质 $(S_k)$ 而言“例外”的 $\xi$ 所成的集合包含由 $x\in P$ 定义的 $P^\vee$ 中的超平面，因而其余维为一（而非至少二）。关于这一方向的一般精确结果，比较下面的（8.5）？
\tnote{首句缺少可确立的完整断言；$f_{k(\xi)}$、$X\subset P^\vee$、点 $X\in X$ 均按底本保留。“不是CM”后面的 $(S_2)$ 句存在文字矛盾，不能当作审定结论。}

\entry{8.4}{命题}
设 $T$ 是 $X$ 的闭子集，并假设 $\codim(T,X)\geq k$。那么，对每个 $\xi\in P^\vee$，都有 $\codim(T_\xi,Y_\xi)\geq k-1$。令 $Z$ 为满足
\[
\codim(T_\xi,Y_\xi)=k-1
\quad\text{（即 }\codim(T_\xi,Y_\xi)<k\text{）}
\]
的 $\xi\in P^\vee$ 所成的集合，则 $Z$ 是 $P^\vee$ 的可构造、无处稠密〔法文：稀疏〕子集，即 $\overline Z$ 在 $P^\vee$ 中的余维至少为一。

它的余维至少为二的充要条件是：对 $X$ 的每个余维等于 $k$ 且满足 $\dim f(T_i)=0$ 的不可约分支 $T_i$，都存在 $X$ 的一个不可约分支 $X_j$，使
\[
\codim(T_i,X_j)=k,\qquad \dim f(X_j)-0
\]
〔或者“闭包”被划去了？〕；即，若 $f$ 拟有限且 $k>0$，则〔原文有不明确的 que (it???)〕$T$ 没有满足 $\dim_xX-k$ 的孤立点。第一个断言立即由下面的引理（8.4.1）(a) 得出；这个引理是对第5节的一项事后补充。
\tnote{此处 $T_i$ 确写为 $X$ 的分支，而不是 $T$ 的分支；当 $k>0$ 时，这与其在 $X$ 中余维等于 $k$ 的要求冲突。减号写法按底本保留，不据猜测填成等号。}

\entry{8.4.1}{引理}
设 $X$ 为局部诺特预概形，$L$ 为 $X$ 上的可逆模，$\phi$ 为 $L$ 的一个截面，$Y=V(\phi)$，$T$ 为 $X$ 的闭子集。假设
\[
\codim(Y,X)\geq k\text{。}
\]
则：
\begin{enumerate}
\item[(a)] $\codim(T\cap Y,Y)\geq k-1$。
\item[(b)] 要有
\[
\codim(T\cap Y,Y)=k-1,
\quad\text{即 }\codim(T\cap Y,Y)<k,
\]
其充要条件是：存在 $T$ 的一个包含于 $Y$ 的不可约分支 $T_i$，使
\[\codim(T_i,X)=k,\]
并且对 $X$ 的每个包含 $T_i$ 且满足
\[
\dim\mathcal O_{X_j,T_i}=\dim\mathcal O_{X,T_i}\quad(=k)
\]
的不可约分支 $X_j$，都有 $X_j\not\subset Y$。
\end{enumerate}
根据 $0_{\mathrm{IV}}$、第IV章中关于维数的一般事实，这个引理的验证是直接的。
\tnote{假设中的 $Y$ 按底本保留。此假设下(b)有明确反例：取 $X=\mathbb A^1$、$L=\mathcal O_X$、$\phi=t$、$Y=\{0\}$、$T=X$、$k=1$。假设成立，且 $\codim(T\cap Y,Y)=0=k-1$；但 $T$ 唯一的不可约分支是 $X$，并不包含于 $Y$，所以所声称的必要条件失败。疑应以 $\codim(T,X)\geq k$ 为假设，但本批不据此替换或声称已经证明修正版。}

在（8.4）的假设下，由（8.4.1）(b) 可以看出哪些超平面 $H_\xi$ 是例外的。排除如下 $\xi\in P^\vee$ 所成的集合 $Z_0$：$T$ 或 $X$ 有一个不可约分支 $R$，满足 $\dim f(R)>0$，且 $R_\xi$ 的“维数过大”。这是一个余维为二的集合，在以下论述中不予计入。例外的 $H_\xi$ 是那些使得存在一个 $T_i$ 满足
\[
\codim(T_i,X)=k,\qquad\dim f(T)=0,\qquad f(T)\ \mathrm{CH}
\]
\ednote{大概是 $H_\xi$。〔底本脚注32。〕}
并且，对 $X$ 的每个满足 $X_j\supset T_i$ 及 $\codim(T_i,X_j)=k$ 的不可约分支 $X_j$，都有 $f(X_j)\not\subset H_\xi$ 的超平面。

给定满足 $\codim(T_i,X)=k$ 的 $T_i$，若存在 $X_j$，使 $\codim(T_i,X_j)=?$〔字迹不清，请询问 Grothendieck〕且 $\dim f(X_j)=0$，则有 $f(X_j)=f(T_i)\not\subset H_\xi$，所以 $\xi$ 相对于 $T_i$ 不是例外的。反之，如果对每个包含 $T_i$ 且满足 $\codim(T_i,X_j)=k$ 的 $X_j$ 都有 $f(X_j)>0$，则对 $\xi\in P^\vee-Z_0$，$\xi$ 相对于 $T_i$ 为例外的充要条件是 $f(T_i)\not\subset H_\xi$；这样的 $\xi$ 所成的集合，是 $P^\vee$ 中一个超平面在 $P-Z_0$ 上的迹。这就证明了（8.4），并且证明了更精确的结果：例外集是一个余维至少为二的集合与若干超平面的并；这些超平面可由以上证明明显地确定。
\tnote{保留了 $f(T)\ \mathrm{CH}$、$f(X_j)>0$、$P-Z_0$ 以及未辨条件，没有补齐。末尾的“不包含”与“超平面”冲突：固定有理点不在 $H_\xi$ 中的参数集合是开集；使超平面包含该点的参数集合才是对偶超平面。}

（恐怕这段写法很松散，或者说很潦草〔英译者措辞说明〕，因为我一直在几何意义下推理，却没有说出来，即取值于代数闭域的点。当然，（8.4）中所陈述的条件确实是几何的，所以可以假设 $k$ 代数闭，并对 $k$-有理点进行论证。）利用（8.4）、（5.7.4）以及（8.2）的末尾，得到：

\entry{8.5}{推论}
假设对所有相伴素循环 $R$，我们有“至多仅仅”〔字迹不清〕\ednote{字迹不清，请询问 A.G.。〔底本脚注33。〕}，并假设 $F$ 满足 $(S_k)$。

要使 $P^\vee$ 中满足 $\phi_\xi$ 为 $F_{k(\xi)}$-正则且 $G_\xi$ 满足 $(S_k)$ 的点所成的（可构造）集合的补集余维至少为二，其充要条件如下：对每个整数 $n\geq0$，以 $Z_n$ 表示满足 $\coprof_x$〔字迹不清〕\ednote{请询问 A.G.。〔底本脚注34。〕}的 $x\in T=\supp F$ 所成的集合。对 $Z_n$ 的每个不可约分支 $Z_{ni}$，若
\[
\codim(Z_{ni},T)=n+k+1,\qquad\dim f(Z_{ni})=0,
\]
则存在 $T$ 的一个包含 $Z_{ni}$ 的不可约分支 $T_j$，使
\[
\codim(Z_{ni},T_j)=n+k+1,\qquad\dim f(T_j)=0\text{。}
\]

当 $f$ 拟有限时，对〔字迹不清，请询问 A.G.〕的每个闭子集 $R$ 都有 $\dim f(R)=\dim R$，所以判据成为：任一 $Z_n$ 中都不存在使 $\dim_zT$（$=\dim F_z$）等于 $n+k+1$ 的孤立点 $z$。当 $F$ 等维且维数为 $d$ 时，若 $d\leq k$，这个条件是空的（实际上我们早已知道，因为这时 $F$ 的 $(S_k)$〔假设〕恰是 Cohen--Macaulay 假设）；若 $d\geq k+1$，其意思是，$T$ 中 $F$ 的余深度 $>d-(k+1)$ 的点所成的集合 $Z_{d-(k+1)}$，也就是说 $F$ 的真深度 $\geq k+1$（尽管事先由 $(S_k)$ 和 $k\leq d$ 只能得到 $F$ 的真深度 $\geq k$）。〔原句在此结束，缺少完整谓语。〕如果不再假设 $F$ 等维，仍可用下面的简单方式表达所需条件：
\tnote{$Z_n$ 的定义和首项假设均有实质缺文，不可据此直接使用判据。即便在维数为 $d$ 的点，以余深度等于维数减深度解释，余深度 $>d-(k+1)$ 推出的是深度 $<k+1$，不是此处印出的 $\geq k+1$。}

\entry{8.6}{}
对每个满足 $\dim F_x\geq k_1$ 的闭点 $x\in\supp F$，都有 $\prof F_x\geq k+1$。

取 $Z=x$，立即可见其充分性。必要性由下述观察得出：对于使 $\phi_\xi$ 为 $X_{k(\xi)}$-正则且 $x\in Y_\xi$ 的每个 $\xi$，都有
\[
\dim G_{\xi,x}=\dim F_x-1,\qquad
\prof G_{\xi,x}=\prof F)_x-1\text{。}
\]
所以，若 $x$ 使上述条件失效，就有 $\prof G_xix\geq k$，但 $\dim G_{\xi,x}\geq k$，这表明 $G_\xi$ 在 $x$ 不满足 $(S_k)$；然而，使 $x\in Y_\xi$ 的 $\xi$ 所成的集合余维为一。〔注意：我隐含地假设了 $k$ 代数闭，而我们立即可以归结到这种情形。〕前述一般判据在（8.6）的情形应当是显然的。
\tnote{首行确印为 $k_1$，不是 $k+1$；$X_{k(\xi)}$、$\prof F)_x$ 及 $G_xix$ 的不协调写法亦保留。若把末段的 $G_xix$ 理解为 $G_{\xi,x}$，则它印出的两条 $\geq k$ 不等式反而保证该点满足 $(S_k)$，与紧接的结论相反。这些疑点不得作为已确立命题传播。}
\backmatter
\chapter*{译校与接续记录}
本批逐段覆盖（8.1）至（8.6），完整保留底本给出的证明文字，没有将“验证直接”扩写为自行创造的证明。排印时仅规范明显的空格、括号闭合及通常的下标排版；具有实质意义的字迹不清、遗漏、异常符号和不等式均保留并说明。第二名审校者对底本第23--25页图像核对了（8.4.1）、（8.4）证明和（8.6）的关键疑点；反例与逻辑方向另行核验。

\textbf{未解决事项：}尚未找到原法文手稿以消解底本缺文；因此本译本不能提供（8.5）缺失的完整假设、$Z_n$ 定义，或（8.6）的审定修正版。译文的忠实性与底本命题的正确性是两个不同问题。

\textbf{下一准确断点：}底本第25页，在（8.6）末句之后，以“We now study the points $y$ of $Y$ that are not smooth”开头的未编号段落。应从那里接译，不得直接跳到（8.7）。本批未合并入任何已有合订本。
\end{document}
